\documentclass[a4paper,10pt]{report}\input{../0.Préambule/Préambule_cours_prof}\begin{document}

\newpage
\setcounter{chapter}{4}
\chapter{Triangle rectangle et cercle circonscrit}

\section{Pour démontrer qu'un point est sur un cercle}

\vspace{-3mm}
\noindent \begin{minipage}{11.5cm}
\begin{prop}
Si un triangle est rectangle alors son cercle circonscrit a pour diamètre son hypoténuse.
\end{prop}

\begin{rmq}
Une autre manière d'énoncer ce théorème : \og Si un triangle est rectangle alors le centre du cercle circonscrit à ce triangle est le milieu de l'hypoténuse. \fg{}
\end{rmq}

\vspace{1cm}
\end{minipage}
\begin{minipage}{5.5cm}
\begin{center}
\definecolor{qqwuqq}{rgb}{0.,0.39215686274509803,0.}
\begin{tikzpicture}[line cap=round,line join=round,>=triangle 45,x=1.0cm,y=1.0cm,scale=1]
\clip(-0.5,-0.1) rectangle (4.5,4.5);
\draw[line width=1.2pt,color=qqwuqq,fill=qqwuqq,fill opacity=0.1] (1.168176829124016,3.642436758981935) -- (1.4150017956609093,3.466883450654597) -- (1.5905551039882473,3.7137084171914903) -- (1.343730137451354,3.8892617255188284) -- cycle; 
\draw [line width=1.2pt] (2.,2.) circle (2.cm);
\draw [line width=1.2pt] (0.,2.)-- (1.343730137451354,3.8892617255188284);
\draw [line width=1.2pt] (1.343730137451354,3.8892617255188284)-- (4.,2.);
\draw [line width=1.2pt] (4.,2.)-- (2.,2.);
\draw [line width=1.2pt] (3.028556614827487,1.9286084629312832) -- (3.028556614827487,2.0713915370687173);
\draw [line width=1.2pt] (2.9714433851725137,1.9286084629312832) -- (2.9714433851725137,2.0713915370687173);
\draw [line width=1.2pt] (2.,2.)-- (0.,2.);
\draw [line width=1.2pt] (1.0285566148274867,1.9286084629312832) -- (1.0285566148274867,2.0713915370687173);
\draw [line width=1.2pt] (0.971443385172513,1.9286084629312832) -- (0.971443385172513,2.0713915370687173);
\begin{scriptsize}
\draw [color=blue] (2.,2.)-- ++(-3.5pt,0 pt) -- ++(7.0pt,0 pt) ++(-3.5pt,-3.5pt) -- ++(0 pt,7.0pt);
\draw[color=blue] (2.23336335347467,1.8039017081317645) node {\large{$M$}};
\draw [color=blue] (4.,2.)-- ++(-3.5pt,0 pt) -- ++(7.0pt,0 pt) ++(-3.5pt,-3.5pt) -- ++(0 pt,7.0pt);
\draw[color=blue] (4.232326391398748,2.389312312095244) node {\large{$B$}};
\draw [color=blue] (0.,2.)-- ++(-3.5pt,0 pt) -- ++(7.0pt,0 pt) ++(-3.5pt,-3.5pt) -- ++(0 pt,7.0pt);
\draw[color=blue] (-0.3510102884128865,2.2608075453715535) node {\large{$A$}};
\draw [color=blue] (1.343730137451354,3.8892617255188284)-- ++(-3.5pt,0 pt) -- ++(7.0pt,0 pt) ++(-3.5pt,-3.5pt) -- ++(0 pt,7.0pt);
\draw[color=blue] (1.2909950641676056,4.331162120364346) node {\large{$C$}};
\end{scriptsize}
\end{tikzpicture}
\end{center}
\end{minipage}

\vspace{-10mm}
\noindent \begin{minipage}{11.5cm}
\begin{ex}
Soit $EFG$ un triangle rectangle en $F$.

\medskip
\noindent Démontre que le point $F$ appartient au cercle de diamètre $[EG]$.
\end{ex}
\end{minipage}

\noindent \begin{minipage}{5.54cm}
\definecolor{qqwuqq}{rgb}{0.,0.39215686274509803,0.}
\begin{tikzpicture}[line cap=round,line join=round,>=triangle 45,x=1.0cm,y=1.0cm]
\clip(-0.5,0.5) rectangle (4.5,4.);
\draw[line width=1.2pt,color=qqwuqq,fill=qqwuqq,fill opacity=0.1] (3.1342410704550265,3.3938334417748757) -- (3.178696114256061,3.0942249017576393) -- (3.478304654273298,3.1386799455586734) -- (3.4338496104722633,3.43828848557591) -- cycle; 
\draw [line width=1.2pt] (3.4338496104722633,3.43828848557591)-- (3.789698861572702,1.0400122614964928);
\draw [line width=1.2pt] (3.789698861572702,1.0400122614964928)-- (2.,2.);
%\draw [line width=1.2pt] (2.8862684238787577,1.4435954107157056) -- (2.953760311672912,1.5694200956631252);
%\draw [line width=1.2pt] (2.8359385498997898,1.470592165833367) -- (2.903430437693944,1.5964168507807865);
\draw [line width=1.2pt] (2.,2.)-- (0.21030113842729775,2.9599877385035076);
%\draw [line width=1.2pt] (1.096569562306056,2.4035831492192137) -- (1.1640614501002102,2.529407834166633);
%\draw [line width=1.2pt] (1.0462396883270881,2.430579904336875) -- (1.1137315761212423,2.5564045892842944);
\draw [line width=1.2pt] (0.21030113842729775,2.9599877385035076)-- (3.4338496104722633,3.43828848557591);
\begin{scriptsize}
\draw [color=blue] (3.789698861572702,1.0400122614964928)-- ++(-3.5pt,0 pt) -- ++(7.0pt,0 pt) ++(-3.5pt,-3.5pt) -- ++(0 pt,7.0pt);
\draw[color=blue] (4.018151780192597,1.4326657153744362) node {\large{$E$}};
\draw [color=blue] (3.4338496104722633,3.43828848557591)-- ++(-3.5pt,0 pt) -- ++(7.0pt,0 pt) ++(-3.5pt,-3.5pt) -- ++(0 pt,7.0pt);
\draw[color=blue] (3.375627946574143,3.874256283124558) node {\large{$F$}};
\draw [color=blue] (0.21030113842729775,2.9599877385035076)-- ++(-3.5pt,0 pt) -- ++(7.0pt,0 pt) ++(-3.5pt,-3.5pt) -- ++(0 pt,7.0pt);
\draw[color=blue] (-0.12255736979299195,3.3316806014023084) node {\large{$G$}};
\end{scriptsize}
\end{tikzpicture}
\end{minipage}
\begin{minipage}{11.3cm}
\noindent\grandscarreaux{\linewidth}{4cm}
\end{minipage}

\vspace{-0.54mm}
\noindent\grandscarreaux{\linewidth}{0.8cm}

\section{Longueur de la médiane}

\vspace{-1.5cm}
\noindent \begin{minipage}{12cm}
\begin{prop}
Si un triangle est rectangle alors la médiane issue du sommet de l'angle droit a pour longueur la moitié de la longueur de l'hypoténuse.
\end{prop}
\end{minipage}
\begin{minipage}{4cm}
\begin{center}
\definecolor{qqwuqq}{rgb}{0.,0.39215686274509803,0.}
\begin{tikzpicture}[line cap=round,line join=round,>=triangle 45,x=1.0cm,y=1.0cm,scale=1]
\clip(-0.5,0.5) rectangle (4.5,5.);
\draw[line width=1.2pt,color=qqwuqq,fill=qqwuqq,fill opacity=0.1] (1.1396911245261034,3.9077844289582626) -- (1.426268072843112,3.8097274003273554) -- (1.5243251014740193,4.096304348644364) -- (1.2377481531570107,4.194361377275271) -- cycle; 
\draw [line width=1.2pt] (1.2377481531570107,4.194361377275271)-- (3.94676024312388,3.2674282180404632);
\draw [line width=1.2pt] (3.94676024312388,3.2674282180404632)-- (2.,2.);
\draw [line width=1.2pt] (3.0362632915671406,2.589465507397848) -- (2.958360309918439,2.7091239033020966);
\draw [line width=1.2pt] (2.9883999332054407,2.5583043147383675) -- (2.910496951556739,2.6779627106426163);
\draw [line width=1.2pt] (2.,2.)-- (0.053239756876120245,0.7325717819595369);
\draw [line width=1.2pt] (1.0895030484432613,1.3220372893573848) -- (1.0116000667945593,1.4416956852616334);
\draw [line width=1.2pt] (1.0416396900815617,1.2908760966979045) -- (0.9637367084328595,1.410534492602153);
\draw [line width=1.2pt] (0.053239756876120245,0.7325717819595369)-- (1.2377481531570107,4.194361377275271);
\draw [line width=1.2pt] (1.2377481531570107,4.194361377275271)-- (2.,2.);
\draw [line width=1.2pt] (1.6769423029896038,3.147582205449224) -- (1.5420650044797182,3.1007300912300013);
\draw [line width=1.2pt] (1.695683148677293,3.0936312860452704) -- (1.5608058501674074,3.046779171826047);
\begin{scriptsize}
\draw [color=blue] (2.,2.)-- ++(-3.5pt,0 pt) -- ++(7.0pt,0 pt) ++(-3.5pt,-3.5pt) -- ++(0 pt,7.0pt);
\draw[color=blue] (2.23336335347467,1.8039017081317645) node {\large{$M$}};
\draw [color=blue] (3.94676024312388,3.2674282180404632)-- ++(-3.5pt,0 pt) -- ++(7.0pt,0 pt) ++(-3.5pt,-3.5pt) -- ++(0 pt,7.0pt);
\draw[color=blue] (4.175213161743774,3.6600816719184066) node {\large{$G$}};
\draw [color=blue] (1.2377481531570107,4.194361377275271)-- ++(-3.5pt,0 pt) -- ++(7.0pt,0 pt) ++(-3.5pt,-3.5pt) -- ++(0 pt,7.0pt);
\draw[color=blue] (1.1767686048576582,4.6310065760529575) node {\large{$F$}};
\draw [color=blue] (0.053239756876120245,0.7325717819595369)-- ++(-3.5pt,0 pt) -- ++(7.0pt,0 pt) ++(-3.5pt,-3.5pt) -- ++(0 pt,7.0pt);
\draw[color=blue] (-0.27961875134416947,1.1042646448583382) node {\large{$E$}};
\end{scriptsize}
\end{tikzpicture}
\end{center}
\end{minipage}

\vspace{-1cm}
\noindent \begin{minipage}{12cm}
\begin{ex}
Le triangle $POT$ est un triangle rectangle en $O$ tel que $TP=8$cm. 

\noindent Le point $S$ est le milieu du segment $[TP]$. 

\medskip
\noindent Quelle est la longueur du segment $[SO]$ ?
\end{ex}
\end{minipage}

\noindent \begin{minipage}{5cm}
\definecolor{qqwuqq}{rgb}{0.,0.39215686274509803,0.}
\begin{tikzpicture}[line cap=round,line join=round,>=triangle 45,x=1.0cm,y=1.0cm]
\clip(-0.5,0.5) rectangle (4.5,4.);
\draw[line width=1.2pt,color=qqwuqq,fill=qqwuqq,fill opacity=0.1] (3.1342410704550265,3.3938334417748757) -- (3.178696114256061,3.0942249017576393) -- (3.478304654273298,3.1386799455586734) -- (3.4338496104722633,3.43828848557591) -- cycle; 
\draw [line width=1.2pt] (3.4338496104722633,3.43828848557591)-- (3.789698861572702,1.0400122614964928);
\draw [line width=1.2pt] (3.789698861572702,1.0400122614964928)-- (2.,2.);
\draw [line width=1.2pt] (2.8862684238787577,1.4435954107157056) -- (2.953760311672912,1.5694200956631252);
\draw [line width=1.2pt] (2.8359385498997898,1.470592165833367) -- (2.903430437693944,1.5964168507807865);
\draw [line width=1.2pt] (2.,2.)-- (0.21030113842729775,2.9599877385035076);
\draw [line width=1.2pt] (1.096569562306056,2.4035831492192137) -- (1.1640614501002102,2.529407834166633);
\draw [line width=1.2pt] (1.0462396883270881,2.430579904336875) -- (1.1137315761212423,2.5564045892842944);
\draw [line width=1.2pt] (0.21030113842729775,2.9599877385035076)-- (3.4338496104722633,3.43828848557591);
\begin{scriptsize}
\draw [color=blue] (2.,2.)-- ++(-3.5pt,0 pt) -- ++(7.0pt,0 pt) ++(-3.5pt,-3.5pt) -- ++(0 pt,7.0pt);
\draw[color=blue] (1.7,1.8039017081317645) node {\large{$S$}};
\draw [color=blue] (3.789698861572702,1.0400122614964928)-- ++(-3.5pt,0 pt) -- ++(7.0pt,0 pt) ++(-3.5pt,-3.5pt) -- ++(0 pt,7.0pt);
\draw[color=blue] (4.018151780192597,1.4326657153744362) node {\large{$T$}};
\draw [color=blue] (3.4338496104722633,3.43828848557591)-- ++(-3.5pt,0 pt) -- ++(7.0pt,0 pt) ++(-3.5pt,-3.5pt) -- ++(0 pt,7.0pt);
\draw[color=blue] (3.375627946574143,3.874256283124558) node {\large{$O$}};
\draw [color=blue] (0.21030113842729775,2.9599877385035076)-- ++(-3.5pt,0 pt) -- ++(7.0pt,0 pt) ++(-3.5pt,-3.5pt) -- ++(0 pt,7.0pt);
\draw[color=blue] (-0.12255736979299195,3.3316806014023084) node {\large{$P$}};
\end{scriptsize}
\end{tikzpicture}
\end{minipage}
\begin{minipage}{11.7cm}
\noindent\grandscarreaux{\linewidth}{4cm}
\end{minipage}

\section{Pour démontrer qu'un triangle est rectangle}

\begin{prop}
Si un triangle est inscrit dans un cercle de diamètre l'un de ses côtés alors il est rectangle et admet ce diamètre pour hypoténuse.
\end{prop}

\begin{ex}
Trace le cercle de diamètre $[SR]$ tel que $SR=7$cm puis place sur ce cercle un point $H$ tel que $RH=4$cm.

\medskip
\noindent Démontre que le triangle $RHS$ est rectangle en $H$.

\medskip
\noindent\grandscarreaux{\linewidth}{8cm}
\end{ex}

\begin{prop}
Si, dans un triangle, la longueur de la médiane relative à un côté est égale à la moitié de la longueur de ce côté alors
ce triangle est rectangle et admet ce côté pour hypoténuse.
\end{prop}

\begin{ex}
$MON$ est un triangle, $U$ est le milieu de $[MN]$ et on a : $MN= 8$cm ; $OU=4$ cm. 

\medskip
\noindent Démontre que le triangle $MON$ est rectangle en $O$.
\end{ex}

\noindent \begin{minipage}{5.54cm}
\begin{tikzpicture}[line cap=round,line join=round,>=triangle 45,x=1.0cm,y=1.0cm]
\clip(-0.5,0.5) rectangle (4.5,4.);
\draw [line width=1.2pt] (3.704029017090242,3.153201758730516)-- (2.,2.);
\draw [line width=1.2pt] (2.9156770938137124,2.5334811964772923) -- (2.835651728100944,2.651730708538331);
\draw [line width=1.2pt] (2.8683772889892976,2.501471050192185) -- (2.788351923276529,2.6197205622532236);
\draw [line width=1.2pt] (2.,2.)-- (0.2959709829097579,0.846798241269484);
\draw [line width=1.2pt] (1.21164807672347,1.3802794377467769) -- (1.1316227110107018,1.4985289498078154);
\draw [line width=1.2pt] (1.1643482718990552,1.3482692914616696) -- (1.0843229061862871,1.4665188035227081);
\draw [line width=1.2pt] (0.17801674959638536,2.956015806706861)-- (2.,2.);
\draw [line width=1.2pt] (1.0968923222717408,2.5544936875729425) -- (1.030550492683201,2.4280588509693346);
\draw [line width=1.2pt] (1.1474662569131842,2.5279569557375274) -- (1.0811244273246443,2.401522119133919);
\begin{scriptsize}
\draw [color=blue] (2.,2.)-- ++(-3.5pt,0 pt) -- ++(7.0pt,0 pt) ++(-3.5pt,-3.5pt) -- ++(0 pt,7.0pt);
\draw[color=blue] (2.23336335347467,1.8039017081317645) node {\large{$U$}};
\draw [color=blue] (3.704029017090242,3.153201758730516)-- ++(-3.5pt,0 pt) -- ++(7.0pt,0 pt) ++(-3.5pt,-3.5pt) -- ++(0 pt,7.0pt);
\draw[color=blue] (3.932481935710136,3.5458552126084597) node {\large{$N$}};
\draw [color=blue] (0.17801674959638536,2.956015806706861)-- ++(-3.5pt,0 pt) -- ++(7.0pt,0 pt) ++(-3.5pt,-3.5pt) -- ++(0 pt,7.0pt);
\draw[color=blue] (0.120173856240646,3.3887938310572823) node {\large{$O$}};
\draw [color=blue] (0.2959709829097579,0.846798241269484)-- ++(-3.5pt,0 pt) -- ++(7.0pt,0 pt) ++(-3.5pt,-3.5pt) -- ++(0 pt,7.0pt);
\draw[color=blue] (-0.036887525310531494,1.2184911041682853) node {\large{$M$}};
\end{scriptsize}
\end{tikzpicture}
\end{minipage}
\begin{minipage}{11.3cm}
\noindent\grandscarreaux{\linewidth}{4cm}
\end{minipage}

\vspace{-0.54mm}
\noindent\grandscarreaux{\linewidth}{3.2cm}

\end{document}